\chapter{Use of Artificial Intelligence Tools}{}
\makeatletter\@mkboth{}{Appendix}\makeatother
\label{appen:ai_usage}

\section{AI Work Declaration}
Claude and Claude Code (Anthropic) were used for the following parts of the project:
\begin{itemize}
    \item \textbf{Literature search}: finding papers to read for the background and literature review (Chapters~\ref{chap:background} and~\ref{chap:literature}).
    \item \textbf{Scripts}: writing scripts to clean the data and run the experiments, such as the scripts that split the dataset, apply the reviewed relabels and run the training of all models.
    \item \textbf{Debugging}: resolving GPU issues encountered during training.
    \item \textbf{Figures}: writing the plotting code for figures in the report, such as the confusion matrices.
    \item \textbf{LaTeX formatting}: formatting of tables, figures and diagrams in the report.
\end{itemize}
All AI-generated code and output was reviewed and tested before use.

\section{Own Work Declaration}
The following parts of the project were done without the use of AI tools:
\begin{itemize}
    \item Formulation of the problem, objectives and scope (Chapter~\ref{chap:introduction}).
    \item Reading of the literature and writing of the background and literature review (Chapters~\ref{chap:background} and~\ref{chap:literature}).
    \item Inspection and manual relabelling of the clarifier images (Chapter~\ref{chap:systemmethods}).
    \item Design of the experiments, including the choice of architectures, data splits, augmentations, class weights and evaluation metrics (Chapter~\ref{chap:experimental_design}).
    \item Running of the experiments, and analysis and interpretation of the results and Grad-CAM heatmaps (Chapter~\ref{chap:results}).
    \item Writing of the text of Chapters~\ref{chap:introduction} to~\ref{chap:conclusion}, and the conclusions drawn from the results.
\end{itemize}
